\begin{ouverture}[geometrie]{15}{Transformations et homothéties}{Partie 4 · Géométrie}
{Symétries · Translation · Rotation · Homothétie · Motifs}
\objectif{Reconnaître et construire une image par transformation.}
\objectif{Utiliser les rapports positifs et négatifs d’homothétie.}
\objectif{Étudier longueurs, aires et volumes.}
\prerequis{Repérage, angles et proportionnalité.}
\end{ouverture}
\pagestyle{otMain}\essentiel\fontsize{11}{15}\selectfont
\essSection{1}{Déplacer ou retourner une figure}
\begin{definition}La symétrie axiale d’axe d associe à A un point A’ tel que d est la médiatrice de [AA’], sauf si A est sur d : alors A’ = A. La symétrie centrale de centre O fait de O le milieu de [AA’].\end{definition}
\begin{definition}Une translation est un glissement défini par une direction, un sens et une distance, par exemple l’envoi d’un point P sur Q. Une rotation de centre O et d’angle donné fait tourner chaque point autour de O, avec un sens précisé.\end{definition}
\begin{propriete}Symétries, translations et rotations conservent longueurs, angles, alignements, aires et parallélisme. Une symétrie axiale retourne l’orientation du motif ; une rotation ou une translation la conserve.\end{propriete}
\begin{savoirfaire}[Construire]Reporter les points remarquables puis les relier dans le même ordre. Une rotation exige même distance au centre et angle orienté correct. Une translation reproduit le même déplacement pour tous les points, sans faire tourner la figure.\end{savoirfaire}
\clearpage
\essSection{2}{Agrandir ou réduire}
\begin{definition}Une homothétie de centre O et de rapport $k\ne0$ envoie A sur A’ avec O, A, A’ alignés et $OA'=|k|\,OA$. Si $k>0$, A et A’ sont sur la même demi-droite issue de O ; si $k<0$, ils sont de part et d’autre de O. O reste fixe.\end{definition}
\begin{propriete}Les longueurs sont multipliées par $|k|$, les aires par $k^2$, les volumes par $|k|^3$. Une droite ne passant pas par O a pour image une droite parallèle. Les angles sont conservés.\end{propriete}
\begin{exemple}Pour $k=-2$, l’image est deux fois plus grande en longueurs, de l’autre côté du centre. Son aire est quatre fois plus grande et son volume huit fois plus grand. Pour $k=-1$, on retrouve la symétrie centrale.\end{exemple}
\begin{savoirfaire}[Frises, pavages, rosaces]Une frise répète un motif par translations suivant une direction. Un pavage recouvre le plan sans trou ni chevauchement. Une rosace organise des rotations autour d’un centre. Rechercher la transformation qui relie les motifs.\end{savoirfaire}\phantomsection\label{cours-pavage}
\label{fig-cours-pavage}\figureGeo{pavage}

\clearpage\rubrique{Découverte}
\begin{blocLivre}{ex-15-1}\begin{exercice}[1]{Une symétrie axiale}
Dans le repère, A(2;1), B(4;1), C(3;3). Construire leurs images par rapport à l’axe vertical des ordonnées.
\label{fig-ex-15-1}\figureGeo{transfo-01}
\end{exercice}\end{blocLivre}
\begin{blocLivre}{ex-15-2}\begin{exercice}[1]{Une symétrie centrale}
O est l’origine et A(3;2). Construire l’image A’ par la symétrie de centre O.
\label{fig-ex-15-2}\figureGeo{transfo-02}
\end{exercice}\end{blocLivre}
\begin{blocLivre}{ex-15-3}\begin{exercice}[1]{Une translation}
La translation envoie P(0;0) sur Q(3;1). Donner l’image du triangle A(1;2), B(3;2), C(2;4).
\label{fig-ex-15-3}\figureGeo{transfo-03}
\end{exercice}\end{blocLivre}
\begin{blocLivre}{ex-15-4}\begin{exercice}[1]{Un quart de tour}
Dans le repère orthonormé, faire tourner A(3;1) de 90 degrés autour de O dans le sens inverse des aiguilles d’une montre.
\label{fig-ex-15-4}\figureGeo{transfo-04}
\end{exercice}\end{blocLivre}
\begin{blocLivre}{ex-15-5}\begin{exercice}[1]{Longueurs conservées}
Un triangle a pour côtés 4 cm, 6 cm et 7 cm. Quel est le périmètre de son image par une rotation de 120 degrés ?
\label{fig-ex-15-5}\figureGeo{transfo-05}
\end{exercice}\end{blocLivre}

\clearpage\rubrique{Entraînement}
\begin{blocLivre}{ex-15-6}\begin{exercice}[2]{Rapport positif}
O est l’origine, A(2;1), B(3;0), C(1;3). Construire l’image par l’homothétie de centre O et rapport 2.
\label{fig-ex-15-6}\figureGeo{transfo-06}
\end{exercice}\end{blocLivre}
\begin{blocLivre}{ex-15-7}\begin{exercice}[2]{Rapport négatif}
Reprendre A(2;1), B(3;0), C(1;3) et construire l’image par l’homothétie de centre O, origine, de rapport $-1{,}5$.
\label{fig-ex-15-7}\figureGeo{transfo-07}
\end{exercice}\end{blocLivre}
\begin{blocLivre}{ex-15-8}\begin{exercice}[2]{Retrouver le rapport}
O, A et A’ sont alignés, OA = 4 cm et OA’ = 10 cm. A et A’ sont de part et d’autre de O. Quel est le rapport ?
\label{fig-ex-15-8}\figureGeo{transfo-08}
\end{exercice}\end{blocLivre}
\begin{blocLivre}{ex-15-9}\begin{exercice}[2]{Périmètre et aire}
Une figure a pour périmètre 18 cm et pour aire 20 cm$^2$. Une homothétie a pour rapport $-3$. Calculer le périmètre et l’aire de l’image.
\label{fig-ex-15-9}\figureGeo{transfo-09}
\end{exercice}\end{blocLivre}
\begin{blocLivre}{ex-15-10}\begin{exercice}[2]{Reconnaître une réduction}
Une maquette est l’image d’un solide par un rapport de 0,4. Le volume initial est 750 cm$^3$. Calculer le volume réduit.
\label{fig-ex-15-10}\figureGeo{transfo-10}
\end{exercice}\end{blocLivre}

\clearpage\rubrique{Problèmes}
\begin{blocLivre}{ex-15-11}\begin{exercice}[3]{Une frise}
Un motif est répété horizontalement tous les 4 cm. Quelle transformation envoie un motif sur le suivant ? Quelle distance entre un point du premier motif et son homologue dans le sixième ?
\label{fig-ex-15-11}\figureGeo{transfo-11}
\end{exercice}\end{blocLivre}
\begin{blocLivre}{ex-15-12}\begin{exercice}[3]{Une rosace}
Une rosace est formée de huit copies régulièrement réparties d’un motif autour de O. Quel angle de rotation envoie une copie sur la suivante ? Après trois pas ?
\label{fig-ex-15-12}\figureGeo{transfo-12}
\end{exercice}\end{blocLivre}
\begin{blocLivre}{ex-15-13}\begin{exercice}[3]{Retrouver un centre}
A(1;1) a pour image A’(3;3) et B(3;1) a pour image B’(7;3) par une homothétie de rapport 2. Déterminer son centre O.
\label{fig-ex-15-13}\figureGeo{transfo-13}
\end{exercice}\end{blocLivre}
\begin{blocLivre}{ex-15-14}\begin{exercice}[3]{Une aire donne-t-elle le signe ?}
Deux triangles homothétiques ont des aires de 12 cm$^2$ et 75 cm$^2$. Donner les rapports possibles du premier vers le second. Peut-on choisir le signe sans connaître leurs positions ?
\label{fig-ex-15-14}\figureGeo{transfo-14}
\end{exercice}\end{blocLivre}
\begin{blocLivre}{ex-15-15}\begin{exercice}[3]{Deux transformations successives}
Une figure subit deux homothéties de même centre O : rapport $-2$, puis rapport $-0{,}5$. Décrire l’image finale d’un point A différent de O.
\label{fig-ex-15-15}\figureGeo{transfo-15}
\end{exercice}\end{blocLivre}

\clearpage
